\subsubsection*{Teambildung}
Die Teambildung für das Projekt erfolgte im Rahmen der Veranstaltung „PRJRMa-Projektstart“ an der WBH.\@ In dieser Veranstaltung wurden diverse Projektideen von verschiedenen Studierenden vorgestellt. Eine dieser Projektideen war die Entwicklung des Coasterbots, welche von Dennis Borutta vorgestellt wurde und zu der es noch kein Team gab. Zur Bearbeitung des Projekts haben sich im Rahmen der Veranstaltung die Studierenden Stefan Etringer und Gereon Such freiwillig bereit erklärt. Damit bestand das Projektteam aus Dennis Borutta, Stefen Etringer und Gereon Such.

\subsubsection*{Rollenverteilung}
Prinzipiell besitzen alle Teammitglieder die Rolle Projektmitarbeiter. Zusätzlich zu dieser Rolle sind noch die Rollen der Projektleitung und des Projektadministators zu besetzen. Die Rolle des Projekleiters hat Dennis Borutta übernommen und führte damit die Meilensteingespräche mit der Projektbetreuung/dem fiktiven Kunden. Die Rolle des Projektadministrators wurde von Stefan Etringer ausgefüllt und pflegte damit die Dokumentenablage im DMS der WBH.\@

Neben diesen Formalen Rollen haben sich die Projektmitglieder thematisch aufgeteilt. So hatte sich Dennis Borutta um die mechanische Konstruktion und den 3D-Druck, Stefan Etringer um die Algorithmik und Simulation und Gereon Such um die Sensorauswahl und Elektronik gekümmert.
